% Draft replacement for the opening paragraphs of the introduction in RealSPR.tex.

In many inverse problems the phase is precisely the part of the signal that the
measurement device does not see.  This phenomenon is already present in the
Pauli problem from quantum mechanics, where one asks how much of a wave function
is encoded by the magnitudes of the function and of its Fourier transform
\cite{pauli1933allgemeinen,ramos-sousa-pauli}; it also appears throughout
optics and signal processing, where detectors typically record intensities
rather than signed or complex amplitudes
\cite{sanz1984phase,MR3069958,eldar2014stability,grohs2020survey}.  The
mathematical abstraction is the phase retrieval problem: recover an unknown
vector or function from phaseless data.  In the real setting the unavoidable
ambiguity is multiplication by a global sign, and the first question is whether
this is the only obstruction.  A substantial finite-dimensional theory has
grown around this question, including injectivity criteria, stability
estimates, convex relaxations, random measurements, and structured models
\cite{bandeira2014saving,balan2013,MR3069958,MR3260258,KS,DB,MR3746047}.
The infinite-dimensional theory is more delicate.  Phase retrieval can hold in
Hilbert spaces \cite{cahill2016infinite}, but the passage from uniqueness to a
uniform inverse estimate is treacherous: stability may fail for continuous
frames \cite{alaifari2017continuous}, and even canonical structured systems
such as Gabor measurements can be severely ill-posed
\cite{alaifari2021gabor,grohs2021stable,grohs20212,steinerberger2020stability}.

This paper is concerned with that second, quantitative layer of the problem.
For applications, qualitative uniqueness is not the end of the story: a
reconstruction principle that collapses under small perturbations is not a
usable measurement scheme.  This is a familiar lesson in modern harmonic
analysis and signal recovery, where sampling, interpolation, and concentration
theorems become robust tools only after one understands their stability
mechanisms.  The function-space point of view of Freeman, Oikhberg, Pineau, and
Taylor \cite{FOTP} makes this issue intrinsic: stable phase retrieval becomes a
geometric property of a subspace of a Banach lattice, and in the Hilbert case it
can be tested on orthogonal pairs.  The present work studies this geometry for
closed spans of independent random variables in $L^2$.  Such spans are a
natural probabilistic analogue of coordinate subspaces, and they sit exactly at
the interface between phase retrieval and the geometry of independent sums.  In
this setting, Calderbank, Daubechies, Freeman, and Freeman
\cite{calderbank2022stable} proposed a sharp characterization in terms of the
$L^1$ sizes of the coordinates; the role of the present paper is to prove the
sufficiency direction up to one exceptional coordinate.  The need for a
quantitative formulation is not cosmetic: examples with merely
H\"older-stable recovery \cite{christ2022examples} show that, in
infinite-dimensional phase retrieval, the modulus of stability is part of the
phenomenon itself.
